\section{Experimental Setup}

Our experimental design separates infrastructure validation from hypothesis testing, enabling us to distinguish technical feasibility (can entropy be extracted?) from performance claims (does entropy outperform max-probability?). This section details the experimental questions, evaluation protocol, and validation criteria.

\subsection{Research Questions}

We structure our evaluation around three experimental questions, each with specific success criteria:

\textbf{RQ1 (Infrastructure):} Can entropy be reliably extracted from frozen LLM forward passes on factual QA tasks?\\
\textbf{Success Criterion:} Extraction rate > 95\% (no technical bottlenecks)\\
\textbf{Measurement:} Fraction of predictions yielding valid entropy values (no NaN, no numerical overflow)

\textbf{RQ2 (Disagreement Patterns):} Do entropy and max-probability disagree in non-trivial proportions?\\
\textbf{Success Criterion:} Q3 quadrant population > 5\% (disagreement cases are not edge cases)\\
\textbf{Measurement:} Fraction of predictions with high max-prob AND high entropy (median splits)

\textbf{RQ3 (Performance Hypothesis):} Does entropy correlate negatively with prediction correctness, indicating it captures error-related uncertainty?\\
\textbf{Success Criterion:} Spearman $\rho < 0$ with $p < 0.05$\\
\textbf{Measurement:} Rank correlation between entropy and binary correctness

RQ1 and RQ2 validate infrastructure independent of model performance. RQ3 tests the hypothesis that entropy signals prediction quality. Critically, RQ3 requires non-zero variance in correctness --- if all predictions are correct or incorrect, correlation is undefined.

\subsection{Dataset Selection and Rationale}

We evaluate on TriviaQA \citep{Joshi2017} unfiltered validation split. TriviaQA consists of 11,313 question-answer pairs sourced from trivia enthusiast websites, covering diverse factual domains (history, geography, entertainment, science). Questions have single-answer targets with multiple acceptable surface forms (e.g., ``Paris'' / ``Paris, France'').

\textbf{Rationale for TriviaQA:}
\begin{enumerate}
\item \textbf{Task Alignment:} Factual QA represents the selective prediction use case --- high-stakes domains (medical QA, legal advice) require models to abstain when uncertain.
\item \textbf{Evaluation Clarity:} Single-answer targets enable unambiguous exact-match evaluation, avoiding subjective correctness judgments.
\item \textbf{Knowledge Intensity:} TriviaQA requires factual knowledge retrieval (not pattern matching), making model capacity directly relevant.
\item \textbf{Established Benchmark:} Widely used in QA literature \citep{Roberts2020,Petroni2019}, facilitating comparison.
\end{enumerate}

We sample 500 examples for proof-of-concept validation, balancing statistical power with compute efficiency. Power analysis ($\alpha=0.05$, power=0.80, medium effect size $r=0.3$) suggests $n=82$ for correlation tests; $n=500$ provides substantial margin.

\subsection{Model Configuration}

\textbf{Planned Model:} Llama-2-7B \citep{Touvron2023}
\begin{itemize}
\item 7 billion parameters, decoder-only transformer
\item 32 layers, 4096 hidden dim, 32 attention heads
\item Trained on 2 trillion tokens
\item Expected TriviaQA accuracy: 10-15\% (zero-shot)
\end{itemize}

\textbf{Actual Model:} GPT-2 \citep{Radford2019}
\begin{itemize}
\item 117 million parameters (1.7\% of planned scale)
\item 12 layers, 768 hidden dim, 12 attention heads
\item Trained on $\sim$40GB text (WebText)
\item Observed TriviaQA accuracy: 0\% (0/500)
\end{itemize}

The model substitution occurred during implementation due to resource constraints (Llama-2-7B requires 14GB GPU memory vs. GPT-2's 500MB). This unintended change became the critical variable exposing our methodological finding.

\subsection{Baseline Methods}

We compare entropy-based uncertainty against two baselines:

\textbf{1. Max-Probability Thresholding (Primary Baseline)}\\
Reject predictions where $\max(p(v))$ falls below threshold $\tau$. \citet{Hendrycks2017} established this as the standard single-forward-pass baseline for OOD detection. Threshold is swept from 0 to 1 to generate coverage-accuracy curves.

\textbf{2. Random Rejection (Sanity Check)}\\
Reject predictions uniformly at random to achieve target coverage. This control verifies that structured uncertainty signals (entropy, max-prob) outperform uninformed rejection.

All methods operate on the same set of predictions from a single frozen forward pass, ensuring fair compute comparison.

\subsection{Evaluation Metrics}

\textbf{Primary Metric: Spearman Rank Correlation ($\rho$)}\\
Measures monotonic association between entropy and binary correctness. Spearman is preferred over Pearson because (1) correctness is binary (non-normal), (2) we test monotonic relationship (not linearity), and (3) Spearman is robust to outliers in entropy. Computed via SciPy's \texttt{spearmanr} function with two-tailed test.

\textbf{Secondary Metrics:}

\begin{enumerate}
\item \textbf{Extraction Rate:} Fraction of predictions yielding valid entropy (no NaN/inf)
\item \textbf{Entropy Range:} $\max(H) - \min(H)$, normalized by $\log|V|$ to assess distribution spread
\item \textbf{Q3 Population:} Fraction in high max-prob, high-entropy quadrant
\item \textbf{Q1 vs Q3 Accuracy Gap:} Accuracy(Q1) - Accuracy(Q3), testing if disagreement cases have lower accuracy
\end{enumerate}

\subsection{Statistical Analysis}

We apply Bonferroni correction for multiple comparisons. With 3 primary tests (extraction rate, correlation, Q3 population), corrected significance threshold is $\alpha/3 = 0.017$. However, extraction rate and Q3 population are directional tests (one-tailed), while correlation is two-tailed, so we report uncorrected p-values and note which survive correction.

For correlation, null hypothesis $H_0$: $\rho = 0$ (no monotonic association). Alternative $H_1$: $\rho < 0$ (negative association, higher entropy predicts lower accuracy). One-tailed test at $\alpha=0.05$.

For quadrant analysis, we use Mann-Whitney U test to compare Q1 vs Q3 accuracy distributions, testing whether median accuracy differs between agreement and disagreement cases.

\subsection{Experimental Protocol}

For each of the 500 TriviaQA examples:

\begin{enumerate}
\item \textbf{Encode Question:} Tokenize question text, truncate to 512 tokens if needed
\item \textbf{Forward Pass:} Run frozen model, extract logits from final token position
\item \textbf{Compute Metrics:}
   \begin{itemize}
   \item Softmax probabilities: $p(v) = \exp(z_v) / \sum \exp(z_{v'})$
   \item Entropy: $H = -\sum p(v) \log p(v)$
   \item Max-probability: $p_{\max} = \max(p(v))$
   \end{itemize}
\item \textbf{Generate Prediction:} Greedy decode (select argmax token)
\item \textbf{Evaluate Correctness:} Exact match against TriviaQA answers (case-insensitive)
\item \textbf{Accumulate:} Store (entropy, max\_prob, correctness) tuples
\end{enumerate}

After processing all examples:

\begin{enumerate}
\setcounter{enumi}{6}
\item \textbf{Statistical Analysis:} Compute Spearman $\rho$, extraction rate, quadrant populations
\item \textbf{Gate Validation:} Check MUST\_WORK criteria (extraction >95\%, $p<0.05$, Q3>5\%)
\item \textbf{Visualization:} Generate figures (scatter, histograms, quadrant plot, gate metrics)
\end{enumerate}

\subsection{Reproducibility}

Code: PyTorch 2.0, Transformers 4.30, SciPy 1.10, Python 3.9\\
Hardware: Single NVIDIA A100 40GB GPU (GPT-2 runs on CPU)\\
Random Seeds: \texttt{torch.manual\_seed(42)}, \texttt{np.random.seed(42)} for sampling\\
Runtime: 15 minutes (GPT-2), estimated 45 minutes (Llama-2-7B)

Full implementation and experiment scripts are available at [repository link]. Dataset is publicly accessible via Hugging Face Datasets (\texttt{trivia\_qa/unfiltered}).

\subsection{Expected Outcomes}

\textbf{Under Valid Conditions (Llama-2-7B with $\sim$10\% accuracy):}
\begin{itemize}
\item Extraction rate: 100\% (no numerical issues expected)
\item Q3 population: 5-15\% (disagreement cases exist)
\item Spearman $\rho$: -0.2 to -0.4 (moderate negative correlation)
\item Interpretation: Infrastructure validated, hypothesis testable (though not necessarily supported)
\end{itemize}

\textbf{Under Invalid Conditions (GPT-2 with 0\% accuracy):}
\begin{itemize}
\item Extraction rate: 100\% (distributions still extractable)
\item Q3 population: 5-15\% (disagreement patterns persist)
\item Spearman $\rho$: NaN (zero-variance correctness)
\item Interpretation: Infrastructure validated, hypothesis untestable
\end{itemize}

Our results fell into the second category, confirming that model capacity gates experiment validity independently of infrastructure success.
